\documentclass[12pt,a4paper]{article}
\usepackage{amsmath,amssymb,amsthm}
\usepackage{mathtools}
\usepackage{hyperref}
\usepackage{geometry}
\usepackage{booktabs}
\usepackage{enumitem}
\usepackage{xcolor}
\usepackage{graphicx}
\graphicspath{{figures/}}
\geometry{margin=2.5cm}
\hypersetup{colorlinks=true,linkcolor=blue,citecolor=blue,urlcolor=blue}

\title{%
  SOLAR/DUNE Significance Analyses\\
  \large Mathematical Derivations: Day-Night, HEP, and Sensitivity%
}
\author{SOLAR/DUNE Analysis}
\date{May 2026}

%% ── Shared notation macros ────────────────────────────────────────────────
%% Rate / count scalars
\newcommand{\Ecal}{\mathcal{E}}          % exposure factor  T M_det
\newcommand{\Mdet}{M_{\mathrm{det}}}     % active detector mass [kt]
\newcommand{\rc}{r_{i}^{c}}             % rate histogram bin i component c
\newcommand{\muc}{\mu_{i}^{c}}           % expected count  = E r^c
\newcommand{\si}{s_{i}}                  % signal count in bin i
\newcommand{\bi}{b_{i}}                  % background count
\newcommand{\nni}{n_{i}}                 % observed/Asimov count

%% Gaussian / Asimov significance
\newcommand{\ZG}{Z^{\mathrm{G}}_{i}}
\newcommand{\ZGerr}{Z^{\mathrm{G,err}}_{i}}
\newcommand{\ZA}{Z_A}
\newcommand{\ZAsig}{Z_A^{\sigma}}

%% Day-night specifics
\newcommand{\fd}{f}                      % day fraction
\newcommand{\fn}{g}                      % night fraction  = 1-f
\newcommand{\Beff}{B^{\mathrm{eff}}_{i}} % effective two-sample background
\newcommand{\seff}{\sigma^{\mathrm{eff}}_{i}} % effective uncertainty
\newcommand{\epstot}{\varepsilon_{\mathrm{tot}}} % combined asymmetry band

%% Profile-likelihood (HEP)
\newcommand{\bhat}{\hat{\beta}}
\newcommand{\srel}{\sigma_{\mathrm{rel}}}
\newcommand{\srelsq}{\sigma_{\mathrm{rel}}^{2}}
\newcommand{\Btot}{B_{\mathrm{tot}}}
\newcommand{\Ntot}{N_{\mathrm{tot}}}
\newcommand{\minexp}{\epsilon_{\min}}

%% Sensitivity specifics
\newcommand{\Tsig}{T^{\mathrm{sig}}_{ij}}
\newcommand{\Tbkg}{T^{\mathrm{bkg}}_{ij}}
\newcommand{\Apred}{A_{\mathrm{pred}}}
\newcommand{\Abkg}{A_{\mathrm{bkg}}}
\newcommand{\spred}{\sigma_{\mathrm{pred}}}
\newcommand{\sbkg}{\sigma_{\mathrm{bkg}}}
\newcommand{\eij}{e_{ij}}
\newcommand{\oij}{o_{ij}}

\begin{document}
\maketitle
\tableofcontents
\newpage

% ===========================================================================
\section{Overview and Scientific Goals}
\label{sec:overview}

This document presents the mathematical derivations underlying the three
significance analyses in the SOLAR/DUNE software framework,
ordered from simplest to most complex:

\begin{enumerate}[leftmargin=*]
  \item \textbf{Day-Night Asymmetry} (\texttt{12DayNight.py}):
    searches for a time-modulated excess of solar neutrinos during nighttime
    relative to daytime, driven by the MSW matter effect inside the Earth.
    Significance is computed from both a Gaussian approximation on the
    rate-difference spectrum and a two-sample Poisson log-likelihood ratio
    (Asimov).

  \item \textbf{HEP Discovery} (\texttt{13HEP.py}):
    searches for the hep solar neutrino flux ($^3\mathrm{He}+p \to
    {}^4\mathrm{He}+e^++\nu_e$) as an absolute excess above background.
    Three significance estimators are computed: Gaussian, Asimov, and a
    profile-likelihood (PL) significance with a single global background
    nuisance, analytically profiled.

  \item \textbf{Sensitivity} (\texttt{14Sensitivity.py}):
    maps the 2D sensitivity contour in the oscillation-parameter plane
    $(\Delta m^2,\,\sin^2\theta_{12})$.
    A Baker-Cousins Poisson deviance with two normalization nuisances
    is minimized numerically over 2D templates in energy and azimuth.
\end{enumerate}

The three analyses share common ingredients (histogram smoothing, thresholds,
adaptive rebinning for HEP) but differ in the statistical complexity of their
model.  Table~\ref{tab:comparison} (Section~\ref{sec:comparison}) summarizes
the key differences.

\medskip
\noindent\textbf{Common notation.}
Throughout, $T$ is the exposure in kt$\cdot$yr,
$\Mdet$ the active detector mass in kt,
and the exposure factor is $\Ecal = T\,\Mdet$.
Index $i$ runs over energy bins; index $j$ over azimuth bins (Sensitivity
only).  Component label $c$ identifies one physical process (signal,
neutron, gamma, radiological, ${}^8\mathrm{B}$).

% ===========================================================================
\section{Common Ingredients}
\label{sec:common}

\subsection{Rate Histograms and Exposure Scaling}
\label{sec:rates}

All three analyses work with per-unit-exposure, per-unit-mass rate
histograms $\rc$ (units: $(\mathrm{yr\cdot kt\cdot MeV})^{-1}$,
integrated over the bin width $\Delta E$).
Expected event counts at exposure $\Ecal$ are:
\begin{equation}
  \muc = \Ecal\,\rc = T\,\Mdet\,\rc.
  \label{eq:rate_to_counts}
\end{equation}
All analyses apply a reconstructed-energy threshold $E_{\mathrm{th}}$
(configured in \texttt{analysis/config.json});
only bins with $E_i \ge E_{\mathrm{th}}$ enter the significance computation.

\subsection{Histogram Smoothing}
\label{sec:smoothing}

Each rate histogram $\rc$ is convolved with a one-dimensional Gaussian
kernel before entering the significance computation:
\begin{equation}
  \tilde{r}_{i}^{c} = \sum_{j} G_\sigma(i-j)\, r_j^{c},
  \qquad
  G_\sigma(k) = \frac{1}{\sqrt{2\pi}\,\sigma}
    \exp\!\left(-\frac{k^2}{2\sigma^2}\right),
  \label{eq:gaussian_smooth}
\end{equation}
implemented via \texttt{scipy.ndimage.gaussian\_filter1d} with
\texttt{mode=`nearest'}.
Whether smoothing is applied to a given component, and the width $\sigma$,
are configured in \texttt{analysis/smoothing.json} under
\texttt{SMOOTHING.ANALYSES.\{ANALYSIS\}.STAGES}.
Separate stages are provided for the fiducial scan and the significance
computation.

\medskip
\noindent\textbf{Non-negativity clipping.}
Gaussian convolution at distribution tails can produce small negative values.
Negative rates are unphysical and cause profile-likelihood divergences
at high exposure (Section~\ref{sec:pl_numerical}).
All smoothed rates are therefore clipped to zero before any significance
computation:
\begin{equation}
  \tilde{r}_{i}^{c} \leftarrow \max\!\left(0,\,\tilde{r}_{i}^{c}\right).
  \label{eq:clip}
\end{equation}

\subsection{Per-Component Background Error Model}
\label{sec:bkg_error}

For each background component $c$ the absolute uncertainty on the raw
rate in bin $i$ combines MC-statistical and systematic contributions in
quadrature:
\begin{equation}
  \sigma_{i}^{c}
  = r_i^{c}\,
    \sqrt{\left(\frac{\delta_i^c}{r_i^c}\right)^{\!2}
          + \left(\srel^c\right)^2},
  \label{eq:bin_error}
\end{equation}
where $\delta_i^c / r_i^c$ is the relative histogram statistical error
and $\srel^c$ is the relative systematic uncertainty
(\texttt{--background\_uncertainty}, typically 2\%).
The total combined background uncertainty in bin $i$ is:
\begin{equation}
  \sigma_i^{\mathrm{bkg}} = \sqrt{\sum_c \bigl(\sigma_i^c\bigr)^2},
  \label{eq:combined_error}
\end{equation}
also convolved with the component-specific smoothing kernel.
For signal components (hep, ${}^8$B) the corresponding quantity uses the
signal systematic $\srel^{\mathrm{sig}} = 30\%$ (HEP) or $4\%$ (Sensitivity).

% ===========================================================================
\section{Day-Night Asymmetry Analysis}
\label{sec:daynight}

\subsection{Physical Observable: MSW-Enhanced Day-Night Effect}
\label{sec:dn_physics}

Solar $\nu_e$ produced in the solar core oscillate in vacuum on their way to
Earth.  During nighttime passages the neutrinos traverse the Earth's interior
and experience additional matter-induced (MSW) oscillations that partially
restore the $\nu_e$ flavour component.  This creates a night-time excess
characterised by the asymmetry:
\begin{equation}
  \Delta_{\mathrm{DN}} \equiv \frac{N_{\mathrm{night}} - N_{\mathrm{day}}}
    {{\tfrac{1}{2}(N_{\mathrm{night}} + N_{\mathrm{day}})}}.
  \label{eq:dn_asymmetry}
\end{equation}
The size of $\Delta_{\mathrm{DN}}$ depends on the Earth electron-density
profile through the MSW potential $V = \sqrt{2}\,G_F\, n_e(r)$
and on the oscillation parameters $(\Delta m^2,\,\sin^2\theta_{12})$.

\subsection{Two-Period Signal and Background Model}
\label{sec:dn_model}

Let $\fd \in (0,1)$ be the fraction of total exposure attributed to
daytime and $\fn = 1 - \fd$ the nighttime fraction
($\fd = 0.5$ by default; the SURF latitude of $44.35^\circ$N gives
$\fd \approx 0.493$ averaged over a full year).
The oscillation-weighted solar rates split into day and night components:
\begin{align}
  r_i^{\mathrm{day}} &= \text{solar rate in daytime sky exposure},\\
  r_i^{\mathrm{night}} &= \text{solar rate in nighttime (Earth-crossing) exposure}.
\end{align}
Isotropic backgrounds (cosmogenic, geological, detector-intrinsic) are
assumed time-uniform and therefore split proportionally between periods.

Event counts observed in day and night periods are:
\begin{align}
  n_i^{\mathrm{night}} &= \Ecal\bigl(\fn\, r_i^{\mathrm{bkg}} + r_i^{\mathrm{night}}\bigr),
  \label{eq:dn_n_night}\\
  n_i^{\mathrm{day}}   &= \Ecal\bigl(\fd\, r_i^{\mathrm{bkg}} + r_i^{\mathrm{day}}\bigr).
  \label{eq:dn_n_day}
\end{align}

The \textbf{asymmetry signal} at scale factor $\xi$ (Section~\ref{sec:dn_band}) is:
\begin{equation}
  \si = \Ecal\cdot\xi\cdot\bigl(r_i^{\mathrm{night}} - r_i^{\mathrm{day}}\bigr).
  \label{eq:dn_signal}
\end{equation}

\subsection{Effective Two-Sample Background}
\label{sec:dn_effective_bkg}

For a two-sample rate-difference test with unequal period lengths $\fd$ and
$\fn$, the optimal significance statistic requires an
\emph{inverse-fraction-weighted} effective background:
\begin{equation}
  \Beff = \frac{n_i^{\mathrm{night}}}{\fn^2}
         + \frac{n_i^{\mathrm{day}}}{\fd^2}.
  \label{eq:dn_Beff}
\end{equation}
In the equal-period limit $\fd=\fn=\tfrac{1}{2}$ and with $\si\ll\Beff$,
this reduces to $\Beff \approx 4\Ecal r_i^{\mathrm{bkg}}$ and the Gaussian
significance becomes the familiar $s_i / \sqrt{b_i}$ form.
Bins where $\Beff = 0$ are zeroed.

\subsection{Gaussian Significance}
\label{sec:dn_gaussian}

Two Gaussian estimators are computed per bin.

\medskip
\noindent\textbf{Simple Gaussian} (no additional background uncertainty):
\begin{equation}
  \ZG = \frac{\si}{\sqrt{\Beff}},
  \label{eq:dn_simple_gaussian}
\end{equation}

\noindent\textbf{Error Gaussian}
(with combined background uncertainty $\seff$, see
Section~\ref{sec:dn_bkg_uncertainty}):
\begin{equation}
  \ZGerr = \frac{\si}{\sqrt{\Beff + \bigl(\seff\bigr)^2}},
  \label{eq:dn_error_gaussian}
\end{equation}
The global significance is the quadrature sum over all bins above threshold:
\begin{equation}
  Z = \sqrt{\sum_{i \ge i_{\mathrm{th}}} \!\left(Z_i\right)^2}.
  \label{eq:dn_global}
\end{equation}
Both raw and Gaussian-smoothed rate versions are used,
producing four output curves per cut per asymmetry-scale value
(\texttt{RawGaussian}, \texttt{RawErrorGaussian},
\texttt{Gaussian}, \texttt{ErrorGaussian})
plus upper/lower band variants (Section~\ref{sec:dn_band}).

\subsection{Background Uncertainty for the Rate-Difference Test}
\label{sec:dn_bkg_uncertainty}

Three independent noise sources are combined in quadrature per period and
per bin, then merged into an effective uncertainty:

\begin{enumerate}[leftmargin=*]
  \item \textbf{Poisson statistical}: $\sqrt{N_{\mathrm{bkg},\,p}}$,
    where $N_{\mathrm{bkg},\,p} = \Ecal\,f_p\,r_i^{\mathrm{bkg}}$ is the
    background count in period $p \in \{\mathrm{day, night}\}$.
  \item \textbf{Normalization systematic}:
    $\srel^{\mathrm{bkg}}\,N_{\mathrm{bkg},\,p}$
    (\texttt{--background\_uncertainty}).
  \item \textbf{Day-fraction systematic}:
    $\delta_f\,\Ecal\,r_i^{\mathrm{bkg}}$
    (\texttt{--day\_fraction\_band}, $\delta_f = 0.01$ by default),
    reflecting uncertainty in the solar zenith angle cut and the run schedule.
\end{enumerate}

Per-period uncertainties:
\begin{equation}
  \sigma_{p,i}
  = \sqrt{N_{\mathrm{bkg},\,p,i}
          + \bigl(\srel^{\mathrm{bkg}}\,N_{\mathrm{bkg},\,p,i}\bigr)^2
          + \bigl(\delta_f\,\Ecal\,r_i^{\mathrm{bkg}}\bigr)^2}.
  \label{eq:dn_sigma_period}
\end{equation}

Combining across periods via inverse-fraction weighting (mirroring
Eq.~\eqref{eq:dn_Beff}):
\begin{equation}
  \seff =
  \sqrt{\left(\frac{\sigma_{\mathrm{night},i}}{\fn}\right)^{\!2}
       +\left(\frac{\sigma_{\mathrm{day},i}}{\fd}\right)^{\!2}},
  \label{eq:dn_sigma_eff}
\end{equation}
set to zero wherever $\Beff = 0$.

\subsection{Two-Sample Poisson Asimov Significance}
\label{sec:dn_asimov}

The Day-Night analysis also computes a two-sample Poisson log-likelihood
ratio (Asimov) as an alternative to the Gaussian estimators.
Under the signal hypothesis (asymmetry scale $\xi$), the Asimov
(expected) observed counts are $n_i^{\mathrm{night}}$ and
$n_i^{\mathrm{day}}$ from Eqs.~\eqref{eq:dn_n_night}--\eqref{eq:dn_n_day}.
The null hypothesis $H_0$ assumes no asymmetry: both periods share the same
per-bin total rate, allocated proportionally to period length:
\begin{equation}
  h_{0,i}^{\mathrm{night}} = \fn\,(n_i^{\mathrm{night}} + n_i^{\mathrm{day}}),
  \qquad
  h_{0,i}^{\mathrm{day}}   = \fd\,(n_i^{\mathrm{night}} + n_i^{\mathrm{day}}).
  \label{eq:dn_h0}
\end{equation}
The two-sample log-likelihood ratio is:
\begin{equation}
  q_i
  = 2\!\left[
    n_i^{\mathrm{night}}\ln\frac{n_i^{\mathrm{night}}}{h_{0,i}^{\mathrm{night}}}
    +
    n_i^{\mathrm{day}}\ln\frac{n_i^{\mathrm{day}}}{h_{0,i}^{\mathrm{day}}}
  \right],
  \label{eq:dn_llr_bin}
\end{equation}
defined as zero whenever any count or null hypothesis count is non-positive.
The global Asimov significance is:
\begin{equation}
  \ZA = \sqrt{\max\!\left(0,\,\sum_{i \ge i_{\mathrm{th}}} q_i\right)}.
  \label{eq:dn_asimov_global}
\end{equation}
This is the natural two-sample extension of the HEP Asimov formula
(Section~\ref{sec:hep_gaussian_asimov}).

\subsection{Asymmetry Uncertainty Band}
\label{sec:dn_band}

Two independent sources of uncertainty on the predicted asymmetry amplitude
are combined in quadrature to define a total band $\epstot$:
\begin{enumerate}[leftmargin=*]
  \item \textbf{Earth density band} $\varepsilon_{\oplus}$
    (\texttt{--earth\_density\_band}, default $0.13$):
    fractional uncertainty from PREM-based oscillation probability calculations.
  \item \textbf{Oscillation parameter band} $\varepsilon_{\mathrm{osc}}$
    (\texttt{--oscillation\_band}, default $0.05$):
    residual uncertainty from PDG $\theta_{12}$ and $\Delta m^2_{21}$ ranges.
\end{enumerate}
\begin{equation}
  \epstot
  = \sqrt{\varepsilon_{\oplus}^2 + \varepsilon_{\mathrm{osc}}^2}
  \approx 0.138 \;\text{(defaults)}.
  \label{eq:dn_band}
\end{equation}
Three scale factors bracket the full predicted range:
\begin{equation}
  \xi \in \{1+\epstot,\;1,\;1-\epstot\}.
  \label{eq:dn_xi}
\end{equation}
The scale acts on the asymmetry signal only (Eq.~\eqref{eq:dn_signal}),
leaving the total event rate and the background unchanged.
This produces three significance curves (upper, nominal, lower) for each
estimator.

\subsection{Results}
\label{sec:dn_results}

Figures~\ref{fig:dn_central_exposure}--\ref{fig:dn_lateral_significance}
show the Day-Night analysis outputs for both HD configurations using the
SolarEnergy estimator at the nominal selection.

\begin{figure}[ht]
\centering
\begin{minipage}{0.48\textwidth}
  \centering
  \includegraphics[width=\textwidth]{dn_central_exposure}
  \caption{centralAPA: exposure needed for $3\sigma$ and $5\sigma$ discovery as a function of day-night asymmetry magnitude.}
  \label{fig:dn_central_exposure}
\end{minipage}\hfill
\begin{minipage}{0.48\textwidth}
  \centering
  \includegraphics[width=\textwidth]{dn_central_significance}
  \caption{centralAPA: significance vs.\ exposure at 30\,kt$\cdot$yr for nominal, upper, and lower asymmetry bands.}
  \label{fig:dn_central_significance}
\end{minipage}
\end{figure}

\begin{figure}[ht]
\centering
\begin{minipage}{0.48\textwidth}
  \centering
  \includegraphics[width=\textwidth]{dn_lateral_exposure}
  \caption{lateralAPA: exposure threshold for $3\sigma$ and $5\sigma$ Day-Night discovery.}
  \label{fig:dn_lateral_exposure}
\end{minipage}\hfill
\begin{minipage}{0.48\textwidth}
  \centering
  \includegraphics[width=\textwidth]{dn_lateral_significance}
  \caption{lateralAPA: significance vs.\ exposure at 30\,kt$\cdot$yr.}
  \label{fig:dn_lateral_significance}
\end{minipage}
\end{figure}

% ===========================================================================
\section{HEP Discovery Analysis}
\label{sec:hep}

\subsection{Signal and Background Model}
\label{sec:hep_signal}

The HEP analysis searches for an absolute rate excess of hep solar neutrinos
above a known background in a one-dimensional energy spectrum above
$E_{\mathrm{th}}$.
Signal (hep) and total background event counts in bin $i$ are:
\begin{align}
  \si &= \Ecal\, r_i^{\mathrm{hep}},
  \label{eq:hep_signal}\\
  \bi &= \Ecal\sum_{c\neq\mathrm{hep}} r_i^c
       = \Ecal\bigl(r_i^{\gamma} + r_i^{n} + r_i^{\mathrm{rad}} + r_i^{{}^8\mathrm{B}}\bigr).
  \label{eq:hep_background}
\end{align}
The Asimov dataset sets observed counts to $\nni = \si + \bi$.

\subsection{Background Error Aggregation}
\label{sec:hep_bkg_error}

Each background component $c$ contributes an absolute uncertainty
$\sigma_i^c$ (Eq.~\eqref{eq:bin_error}).
The combined error and its smoothed counterpart enter the significance
computation as:
\begin{equation}
  \sigma_i^{\mathrm{bkg}} = \sqrt{\sum_{c \neq \mathrm{hep}} \bigl(\sigma_i^c\bigr)^2}.
  \label{eq:hep_combined_err}
\end{equation}

\subsection{Adaptive Tail Rebinning}
\label{sec:rebin}

\subsubsection{Detection Threshold}
\label{sec:detection_threshold}

A bin (or merged group) is declared \emph{detectable} if the expected
(conservatively shifted) signal exceeds:
\begin{equation}
  T_{\mathrm{det}} = \max\!\left(T_{\mathrm{min}},\;-\ln(1-p_{\mathrm{min}})\right),
  \label{eq:detection_threshold}
\end{equation}
with defaults $T_{\mathrm{min}}=1.0$ and $p_{\mathrm{min}}=1-e^{-1}$,
giving $T_{\mathrm{det}}=1.0$.
The Poisson interpretation: $P(\ge 1;\lambda)\ge p \Leftrightarrow \lambda\ge -\ln(1-p)$.

The conservative (downward-fluctuated) detectable signal is:
\begin{equation}
  s_i^{\mathrm{det}} = \si\,(1 - d\cdot\srel^{\mathrm{sig}}),
  \quad d \in \{2.9,\;3.0,\;3.1\},
  \label{eq:detectable_signal}
\end{equation}
where the three values of $d$ correspond to the upper (+), nominal, and
lower ($-$) error bands respectively.

\subsubsection{Greedy Algorithm and Monotonicity Enforcement}

Bins are merged from the high-energy tail inward until each group exceeds
$T_{\mathrm{det}}$.  Groups below $T_{\mathrm{det}}$ are zeroed out.
To guarantee a non-decreasing significance vs.\ exposure curve, if the
current step's optimal binning yields lower Asimov significance than the
previous step the previous binning is retained:
\begin{equation}
  \mathcal{B}(t) =
  \begin{cases}
    \mathcal{B}^*(t) & Z_A(\mathcal{B}^*(t),t) \ge Z_A(\mathcal{B}(t{-}1),t{-}1),\\
    \mathcal{B}(t-1) & \text{otherwise.}
  \end{cases}
  \label{eq:monotone_rebin}
\end{equation}

\subsection{Gaussian and Asimov Significance}
\label{sec:hep_gaussian_asimov}

With $\si$ and $\bi$ from Eqs.~\eqref{eq:hep_signal}--\eqref{eq:hep_background}
and combined background uncertainty $\sigma_i^{\mathrm{bkg}}$,
the per-bin Gaussian estimator is:
\begin{equation}
  \ZG = \frac{\si}{\sqrt{\bi + \bigl(\sigma_i^{\mathrm{bkg}}\bigr)^2}},
  \label{eq:hep_gaussian}
\end{equation}
and the global significance is Eq.~\eqref{eq:dn_global}.

The Asimov significance \cite{Cowan2010}, with and without background uncertainty, is:
\begin{align}
  Z_A &= \sqrt{2\left[(\si+\bi)\ln\!\left(1+\frac{\si}{\bi}\right) - \si\right]},
  \label{eq:asimov_simple}\\
  Z_A^{\sigma} &= \sqrt{2\left[
    (\si+\bi)\ln\!\frac{(\si+\bi)(\bi+(\sigma_i^{\mathrm{bkg}})^2)}
                       {\bi^2+(\si+\bi)(\sigma_i^{\mathrm{bkg}})^2}
    -\frac{\bi^2}{(\sigma_i^{\mathrm{bkg}})^2}
     \ln\!\left(1+\frac{(\sigma_i^{\mathrm{bkg}})^2 \si}{\bi(\bi+(\sigma_i^{\mathrm{bkg}})^2)}\right)
  \right]}.
  \label{eq:asimov_full}
\end{align}
These are computed both with and without adaptive tail rebinning.
The profile-likelihood (Section~\ref{sec:pl}) is the primary metric.

\subsection{Profile-Likelihood Significance}
\label{sec:pl}

\subsubsection{Test Statistic}
\label{sec:pl_model}

The profile-likelihood significance follows Cowan et al.\ \cite{Cowan2010}.
The test statistic for the null hypothesis $\mu=0$ (no HEP signal) is:
\begin{equation}
  q_0 = -2\ln\!\frac{\mathcal{L}(0,\,\bhat)}{\mathcal{L}(\hat{s}+b,\,1)},
  \label{eq:q0}
\end{equation}
evaluated at the Asimov point $\nni = \si + \bi$, with $\bhat$ the
background scale factor profiled under the null.
The median discovery significance is $Z = \sqrt{q_0}$.

\subsubsection{Global Background Normalization Nuisance}
\label{sec:pl_nuisance}

A single global background scale factor $\beta$ — fully correlated across
all bins — is constrained by a Gaussian penalty:
\begin{equation}
  \mathcal{L}(0,\beta)
  = \prod_{i=1}^{K}
      \frac{(\beta \bi)^{\nni} e^{-\beta \bi}}{\nni!}
    \cdot\exp\!\left(-\frac{(\beta-1)^2}{2\srelsq}\right),
  \label{eq:likelihood}
\end{equation}
where $\srel$ is the relative background uncertainty
(\texttt{--background\_uncertainty}, typically 2\%).
A global (fully correlated) $\beta$ avoids the $K$-parameter absorbing
plateau that per-bin nuisances produce; the plateau ends at
$T \sim 1/(B_{\mathrm{tot}}\srelsq)$ (typically sub-year),
giving a physically correct PL curve.

\subsubsection{Profiling \texorpdfstring{$\bhat$}{beta-hat}}
\label{sec:betahat}

The stationarity condition
$\partial\ln\mathcal{L}/\partial\beta\big|_{\bhat}=0$
at the Asimov point $\nni=\si+\bi$ gives:
\begin{equation}
  \Ntot = \bhat\Btot + \frac{\bhat-1}{\srelsq},
  \quad
  \Ntot = \sum_i \nni,\;\Btot = \sum_i \bi.
  \label{eq:stationarity}
\end{equation}
Because $\beta$ is \emph{global}, the per-bin denominators $\beta\bi$
factorize and the stationarity condition is a scalar equation.
Rearranging:
\begin{equation}
  \bhat^2 + \left(\Btot\srelsq-1\right)\bhat - \Ntot\srelsq = 0,
  \label{eq:beta_quadratic}
\end{equation}
with analytic positive root:
\begin{equation}
  \bhat = \frac{-(\Btot\srelsq-1)
    + \sqrt{(\Btot\srelsq-1)^2 + 4\Ntot\srelsq}}{2}.
  \label{eq:beta_solution}
\end{equation}
When $\Btot=0$ or $\srel=0$, $\bhat=1$ (no constraint active).

\subsubsection{Log-Likelihood Ratio}
\label{sec:llr}

Substituting $\nni = \si+\bi$ into Eq.~\eqref{eq:q0}:
\begin{equation}
  q_0 = 2\sum_{i=1}^{K}
    \left[\nni\ln\!\frac{\nni}{\bhat \bi} - (\nni-\bhat \bi)\right]
    + \left(\frac{\bhat-1}{\srel}\right)^{\!2}.
  \label{eq:q0_expanded}
\end{equation}
Each per-bin term is the Baker-Cousins Poisson deviance \cite{BakerCousins1984}
between the Asimov observation $\nni$ and the null expectation $\bhat\bi$.
This expression — a sum of Poisson deviances with a global nuisance penalty
— is the structural prototype for the Sensitivity objective
(Section~\ref{sec:sens_objective}), extended to 2D with two nuisances.

\subsubsection{Numerical Stability}
\label{sec:pl_numerical}

When $S^2/B \ll 1$, naive subtraction of two $\mathcal{O}(N\log N)$
log-likelihoods loses precision.  The stable per-bin form is:
\begin{equation}
  \Delta\ell_i = \nni\ln\!\frac{\nni}{\mu_i^{\mathrm{null}}}
    - \bigl(\nni - \mu_i^{\mathrm{null}}\bigr),
  \qquad
  \mu_i^{\mathrm{null}} = \bhat \bi,
  \label{eq:stable_llr}
\end{equation}
with a floor $\minexp=10^{-12}$ replacing zero denominators.

\medskip
\noindent\textbf{Negative-rate blowup.}
If a smoothed rate is negative, $\bhat\bi<0$, the floor gives
$\Delta\ell_i \approx \nni\ln(\nni/\minexp) \propto T$,
causing a spurious superlinear spike.
Clipping (Eq.~\eqref{eq:clip}) eliminates this path entirely.

\subsection{Barlow-Beeston MC Mask}
\label{sec:bb_mask}

Bins with summed background MC count below $N_{\mathrm{MC}}^{\min}$
(default 1.0) lack a reliable background model and are excluded by a
static binary mask \cite{BarlowBeeston1993}:
\begin{equation}
  \mathcal{M}_i = \mathbf{1}\!\left[N_i^{\mathrm{MC}} \ge N_{\mathrm{MC}}^{\min}\right],
  \qquad
  \si \leftarrow \mathcal{M}_i \si,\quad \bi \leftarrow \mathcal{M}_i \bi.
  \label{eq:bb_mask_hep}
\end{equation}
The mask is static (MC counts do not change with exposure), so it introduces
no discrete transitions in the significance curve.
The same mask is applied to both raw and smoothed histograms.

\subsection{Profile-Likelihood Signal Bands and Post-Processing}
\label{sec:pl_bands}

Three signal normalizations are evaluated to produce $\pm1\sigma_s$ bands,
using the same scale factor $d \in \{+1, 0, -1\}$ applied to the signal:
\begin{equation}
  s_i^{(d)} = \si\,(1 + d\cdot\srel^{\mathrm{sig}}).
  \label{eq:signal_variation}
\end{equation}
Background is never shifted; $\bhat$ and its quadratic solution are
unaffected.

The raw PL significance curve $Z(T)$ is post-processed by two steps when
\texttt{pl\_isotonic} is enabled (controlled by \texttt{analysis/config.json}
under \texttt{WORKFLOW.HEP.pl\_isotonic}):
\begin{enumerate}[leftmargin=*]
  \item \textbf{Gaussian kernel smoothing} with $\sigma_{\mathrm{PL}}=6$
    exposure-grid index units
    (\texttt{scipy.ndimage.gaussian\_filter1d}), which suppresses numerical
    oscillations from the PL solver at low signal-to-background ratio.
  \item \textbf{Isotonic regression (PAVA)} \cite{Robertson1988}:
    \begin{equation}
      \min_{\{z_k\}}\sum_k(z_k-\tilde{Z}(T_k))^2
      \quad\text{s.t.}\quad z_k\le z_{k+1},
      \label{eq:pava}
    \end{equation}
    via \texttt{sklearn.isotonic.IsotonicRegression},
    which enforces a non-decreasing significance curve while minimising the
    $\ell^2$ deviation from the computed values.
\end{enumerate}
Pipeline: $Z \to \max(0,Z) \xrightarrow{\mathrm{Gauss}} \tilde{Z}
          \xrightarrow{\mathrm{PAVA}} Z'$.
After isotonic regression, upper/lower bands are enforced to remain
above/below the nominal curve respectively:
$Z'^{(+1)}\leftarrow\max(Z'^{(+1)}, Z'^{(0)})$ and
$Z'^{(-1)}\leftarrow\min(Z'^{(-1)}, Z'^{(0)})$.
Pre-PAVA curves are saved separately for diagnostic purposes.

\subsection{Spike Detection and Best-Cut Selection}
\label{sec:spike}

A PL curve is declared \emph{spiked} if any consecutive step in the
pre-isotonic columns exceeds $\Delta_{\max}$
(\texttt{--max\_pl\_jump}, default $1\,\sigma$):
\begin{equation}
  \text{spiked} \;\Leftrightarrow\;
  \max_k\!\left[Z_{\mathrm{pre}}(T_{k+1})-Z_{\mathrm{pre}}(T_k)\right]>\Delta_{\max}.
  \label{eq:spike_criterion}
\end{equation}
The best cut maximises the reference significance over non-spiked cuts:
\begin{equation}
  (N^*,N_{\mathrm{op}}^*,N_{\mathrm{adj}}^*) =
  \arg\max_{\text{not spiked}} Z_{\mathrm{ref}}(T_{\max}).
  \label{eq:best_cut}
\end{equation}

\subsection{Results}
\label{sec:hep_results}

Figures~\ref{fig:hep_central_exposure}--\ref{fig:hep_lateral_adaptive}
show HEP analysis outputs for both HD configurations (SolarEnergy, nominal).
Each row presents: left, Gaussian/Asimov exposure comparison; right,
profile-likelihood significance with adaptive rebinning.

\begin{figure}[ht]
\centering
\begin{minipage}{0.48\textwidth}
  \centering
  \includegraphics[width=\textwidth]{hep_central_exposure}
  \caption{centralAPA: Gaussian and Asimov exposure needed to reach $3\sigma$ and $5\sigma$ hep discovery.}
  \label{fig:hep_central_exposure}
\end{minipage}\hfill
\begin{minipage}{0.48\textwidth}
  \centering
  \includegraphics[width=\textwidth]{hep_central_pl_exposure}
  \caption{centralAPA: profile-likelihood significance vs.\ exposure.}
  \label{fig:hep_central_pl_exposure}
\end{minipage}
\end{figure}

\begin{figure}[ht]
\centering
\begin{minipage}{0.48\textwidth}
  \centering
  \includegraphics[width=\textwidth]{hep_central_pl_adaptive}
  \caption{centralAPA: PL significance with and without adaptive rebinning.}
  \label{fig:hep_central_pl_adaptive}
\end{minipage}\hfill
\begin{minipage}{0.48\textwidth}
  \centering
  \includegraphics[width=\textwidth]{hep_central_asimov_adaptive}
  \caption{centralAPA: Asimov significance with and without adaptive rebinning.}
  \label{fig:hep_central_asimov_adaptive}
\end{minipage}
\end{figure}

\begin{figure}[ht]
\centering
\begin{minipage}{0.48\textwidth}
  \centering
  \includegraphics[width=\textwidth]{hep_lateral_exposure}
  \caption{lateralAPA: Gaussian and Asimov exposure threshold for hep discovery.}
  \label{fig:hep_lateral_exposure}
\end{minipage}\hfill
\begin{minipage}{0.48\textwidth}
  \centering
  \includegraphics[width=\textwidth]{hep_lateral_pl_exposure}
  \caption{lateralAPA: profile-likelihood significance vs.\ exposure.}
  \label{fig:hep_lateral_pl_exposure}
\end{minipage}
\end{figure}

\begin{figure}[ht]
\centering
\begin{minipage}{0.48\textwidth}
  \centering
  \includegraphics[width=\textwidth]{hep_lateral_pl_adaptive}
  \caption{lateralAPA: PL significance with and without adaptive rebinning.}
  \label{fig:hep_lateral_pl_adaptive}
\end{minipage}\hfill
\begin{minipage}{0.48\textwidth}
  \centering
  \includegraphics[width=\textwidth]{hep_lateral_asimov_adaptive}
  \caption{lateralAPA: Asimov significance with and without adaptive rebinning.}
  \label{fig:hep_lateral_adaptive}
\end{minipage}
\end{figure}

% ===========================================================================
\section{Sensitivity Analysis}
\label{sec:sensitivity}

\subsection{Overview}
\label{sec:sens_overview}

The Sensitivity analysis maps how well the detector can distinguish between
the solar ($\Delta m^2_\odot$) and reactor ($\Delta m^2_{\mathrm{react}}$)
best-fit oscillation hypotheses.
Rather than computing a single discovery significance it evaluates a
$\chi^2$ surface over the oscillation-parameter grid
$(\Delta m^2,\,\sin^2\theta_{13},\,\sin^2\theta_{12})$.

Key differences from HEP:
\begin{itemize}[leftmargin=*]
  \item Histograms are \textbf{two-dimensional} (energy $\times$ azimuth $\cos\eta$).
  \item \textbf{Two} normalization nuisances are floated: $\Apred$ (signal)
    and $\Abkg$ (background).
  \item The nuisance stationarity conditions have
    \textbf{no closed-form solution} — unlike HEP's quadratic for $\bhat$
    — requiring full 2D numerical minimization.
  \item The goal is a $\chi^2$ \emph{map}, not a single significance.
\end{itemize}

\subsection{Two-Dimensional Template Construction}
\label{sec:sens_templates}

Signal and background are represented as 2D arrays with axes
energy ($j=1,\ldots,J$) and azimuth ($i=1,\ldots,I$).

\medskip
\noindent\textbf{Signal template.}
For oscillation parameters $\vec{\theta} = (\Delta m^2,\sin^2\theta_{13},\sin^2\theta_{12})$,
the signal template at exposure $\Ecal$ is:
\begin{equation}
  \Tsig(\vec{\theta}) = \Ecal\cdot\Mdet\cdot\bigl[P(\vec{\theta})\,H\bigr]_{ij},
  \label{eq:signal_template}
\end{equation}
where $P(\vec{\theta})$ is the oscillation-probability matrix
(azimuth $\times$ neutrino flavour) and $H$ is the detector energy-response
matrix (neutrino energy $\times$ reconstructed energy),
implemented in \texttt{14SensitivitySignalTemplate.py}.

\medskip
\noindent\textbf{Background template.}
The background template $\Tbkg$ is independent of oscillation parameters
and is constructed from detector simulations in
\texttt{14SensitivityBackgroundTemplate.py}.

\subsection{Asimov Dataset Construction}
\label{sec:sens_asimov}

For each scan point $\vec{\theta}_k$ the Asimov observed dataset is:
\begin{equation}
  \oij(\vec{\theta}_k) = \Tsig(\vec{\theta}_k) + \Tbkg.
  \label{eq:sens_asimov}
\end{equation}
Two \emph{reference} templates fixed at the solar and reactor best-fit
points:
\begin{equation}
  p^{\mathrm{solar}}_{ij} = \Tsig(\vec{\theta}_{\odot}),
  \qquad
  p^{\mathrm{react}}_{ij} = \Tsig(\vec{\theta}_{\mathrm{react}}).
  \label{eq:reference_templates}
\end{equation}

\subsection{Objective Function: Baker-Cousins Poisson Deviance}
\label{sec:sens_objective}

The fit minimizes a Baker-Cousins Poisson deviance \cite{BakerCousins1984}
with two free normalization nuisances:
\begin{equation}
  \chi^2(\Apred,\Abkg)
  = 2\sum_{i,j}\Delta\ell_{ij}
    + \left(\frac{\Apred}{\spred}\right)^{\!2}
    + \left(\frac{\Abkg}{\sbkg}\right)^{\!2},
  \label{eq:sens_chisq}
\end{equation}
where the expected model is:
\begin{equation}
  \eij = (1+\Abkg)\,\Tbkg + (1+\Apred)\,p_{ij},
  \label{eq:sens_expected}
\end{equation}
and the per-bin Poisson deviance is:
\begin{equation}
  \Delta\ell_{ij} =
  \begin{cases}
    \eij - \oij + \oij\ln(\oij/\eij) & \oij>0,\; \eij>0,\\
    \eij & \oij=0,\; \eij>0,\\
    0 & \eij=0 \;\text{(or masked)}.
  \end{cases}
  \label{eq:sens_deviance}
\end{equation}
Each term $\Delta\ell_{ij}\ge0$ by the Gibbs inequality; the sum is zero
if and only if $\eij=\oij$ for all bins.

\medskip
\noindent\textbf{Structural connection to HEP.}
Under $H_0$ (no signal, $p_{ij}=0$), Eq.~\eqref{eq:sens_chisq} reduces to
the HEP log-likelihood ratio Eq.~\eqref{eq:q0_expanded} with one nuisance
per normalisation parameter.  The Sensitivity objective is therefore the
direct 2D two-nuisance generalization of the HEP test statistic.

\subsection{Nuisance Parameters: Constraints and Non-Analyticity}
\label{sec:sens_nuisances}

The penalty terms in Eq.~\eqref{eq:sens_chisq} are Gaussian constraints
identical in form to the HEP penalty $[(\bhat-1)/\srel]^2$, but with
separate widths: $\spred = 4\%$ (signal flux uncertainty) and
$\sbkg = 2\%$ (background uncertainty).
These are configured via \texttt{analysis/config.json} under
\texttt{ANALYSIS\_UNCERTAINTIES.SENSITIVITY}.

\medskip
\noindent\textbf{Why no analytic solution exists.}
In HEP, $e_i = \beta\bi$ under $H_0$:
the single global $\beta$ factors out and the stationarity condition
collapses to the scalar quadratic Eq.~\eqref{eq:beta_quadratic}.

In Sensitivity, $\eij=(1+\Abkg)\Tbkg+(1+\Apred)p_{ij}$ mixes $\Tbkg$
and $p_{ij}$ in every bin.  The two stationarity conditions are:
\begin{align}
  \sum_{ij} p_{ij}\!\left(1-\frac{\oij}{\eij}\right)
    &= -\frac{\Apred}{\spred^2},
  \label{eq:sens_stat_pred}\\
  \sum_{ij} \Tbkg\!\left(1-\frac{\oij}{\eij}\right)
    &= -\frac{\Abkg}{\sbkg^2},
  \label{eq:sens_stat_bkg}
\end{align}
which are jointly non-linear in $(\Apred,\Abkg)$ because $\eij$ in
the denominators depends on both unknowns.  Full 2D minimization is
required.

\subsection{Barlow-Beeston Mask for 2D Templates}
\label{sec:sens_bb}

Bins where the background template is zero are excluded:
\begin{equation}
  \mathcal{M}_{ij} = \mathbf{1}\!\left[\Tbkg > 0\right],
  \label{eq:sens_bb_mask}
\end{equation}
passed as \texttt{bb\_mask} to \texttt{Sensitivity\_Fitter}.
This prevents spurious large deviance from signal-only bins
lacking a reliable background model,
following the same spirit as HEP's Barlow-Beeston mask
(Section~\ref{sec:bb_mask}).

\subsection{Minimization: scipy L-BFGS-B}
\label{sec:sens_minimization}

The objective Eq.~\eqref{eq:sens_chisq} is minimized over
$(\Apred,\Abkg)$ using the L-BFGS-B algorithm
\cite{ByrdEtAl1995,ZhuEtAl1997}:
\begin{equation}
  (\hat{A}_{\mathrm{pred}},\hat{A}_{\mathrm{bkg}})
  = \arg\min_{\Apred,\Abkg}\chi^2(\Apred,\Abkg),
  \label{eq:sens_minimize}
\end{equation}
with box constraints:
\begin{equation}
  |\Apred| \le 10\,\spred,
  \qquad
  |\Abkg| \le 10\,\sbkg.
  \label{eq:sens_bounds}
\end{equation}
L-BFGS-B uses gradient information and is well-suited to smooth,
convex objectives with box constraints.
Implemented in \texttt{lib/lib\_root.py:Sensitivity\_Fitter.Fit}.

\subsection{Oscillation Grid Scan and Cut Quality Score}
\label{sec:sens_scan}

For each analysis cut $(N_{\mathrm{hits}},N_{\mathrm{ophits}},N_{\mathrm{adjcl}})$
and each grid point $\vec{\theta}_k$:
\begin{enumerate}[leftmargin=*]
  \item Construct the Asimov dataset $\oij(\vec{\theta}_k)$.
  \item Minimize $\chi^2$ against $p^{\mathrm{solar}}$:
    obtain $\chi^2_\odot(\vec{\theta}_k)$.
  \item Minimize $\chi^2$ against $p^{\mathrm{react}}$:
    obtain $\chi^2_{\mathrm{react}}(\vec{\theta}_k)$.
\end{enumerate}

The \textbf{cut quality score} is the average wrong-hypothesis $\chi^2$
(larger $\Rightarrow$ better discrimination):
\begin{equation}
  \mathrm{Score}(\mathrm{cut}) =
    \tfrac{1}{2}\!\left[
      \chi^2_\odot(\vec{\theta}_{\mathrm{react}})
      + \chi^2_{\mathrm{react}}(\vec{\theta}_\odot)
    \right].
  \label{eq:sens_score}
\end{equation}
The resulting $\chi^2$ surface over the oscillation grid gives the
sensitivity contours.

\subsection{Template Construction Pipeline}
\label{sec:sens_pipeline}

The 2D templates $T^{\mathrm{sig}}$ and $T^{\mathrm{bkg}}$ are assembled
from three physical inputs via a two-stage convolution.
Figures~\ref{fig:sens_conv_nadir}--\ref{fig:sens_conv_oscillogram} show the
nadir time-fraction $p(\cos\eta)$ (earth geometry, fixed by DUNE latitude
$44.35^\circ$) and the solar survival probability
$P(\nu_e\!\to\!\nu_e;\,E_\nu^{\mathrm{true}},\cos\eta)$ computed at the
best-fit oscillation parameters.
Figures~\ref{fig:sens_conv_background1d}--\ref{fig:sens_conv_signal1d}
present the 1D background spectrum $b(E_{\mathrm{reco}})$ from the selected
cuts and the 1D marginal signal spectrum $s(E_{\mathrm{reco}})$ obtained by
summing the oscillated signal over nadir bins.
Figures~\ref{fig:sens_conv_background_template}--\ref{fig:sens_conv_signal_template}
show the fully assembled 2D templates:
$T^{\mathrm{bkg}}$ is formed by weighting $b(E_{\mathrm{reco}})$ by
$p(\cos\eta)$ via \texttt{\_project\_1d\_to\_2d};
$T^{\mathrm{sig}}$ is the result of the full convolution
$T^{\mathrm{sig}} = P_{\nu_e\to\nu_e}\cdot h^T$,
where $h$ is the $(E_{\mathrm{reco}}\times E_{\mathrm{true}})$ smearing matrix.

\begin{figure}[ht]
\centering
\begin{minipage}{0.48\textwidth}
  \centering
  \includegraphics[width=\textwidth]{sens_conv_nadir}
  \caption{Nadir time-fraction $p(\cos\eta)$ at DUNE latitude ($44.35^\circ$N),
    derived from solar ephemeris.  This distribution weights each nadir strip in
    the 2D background template.}
  \label{fig:sens_conv_nadir}
\end{minipage}\hfill
\begin{minipage}{0.48\textwidth}
  \centering
  \includegraphics[width=\textwidth]{sens_conv_oscillogram}
  \caption{Solar-neutrino oscillogram $P(\nu_e\!\to\!\nu_e;\,E_\nu,\cos\eta)$
    at best-fit parameters
    ($\Delta m^2_{21}=6\times10^{-5}$\,eV$^2$, $\sin^2\theta_{12}=0.303$).
    Each row is the survival-probability spectrum for one nadir bin.}
  \label{fig:sens_conv_oscillogram}
\end{minipage}
\end{figure}

\begin{figure}[ht]
\centering
\begin{minipage}{0.48\textwidth}
  \centering
  \includegraphics[width=\textwidth]{sens_conv_background1d}
  \caption{1D background spectrum $b(E_{\mathrm{reco}})$ after optimal cut
    selection, showing individual components and the smoothed total used as
    input to the 2D background template.}
  \label{fig:sens_conv_background1d}
\end{minipage}\hfill
\begin{minipage}{0.48\textwidth}
  \centering
  \includegraphics[width=\textwidth]{sens_conv_signal1d}
  \caption{1D marginal signal spectrum $s(E_{\mathrm{reco}})=\sum_\eta T^{\mathrm{sig}}$
    at solar best-fit parameters and $30\,\mathrm{kt}\!\cdot\!\mathrm{yr}$
    exposure.}
  \label{fig:sens_conv_signal1d}
\end{minipage}
\end{figure}

\begin{figure}[ht]
\centering
\begin{minipage}{0.48\textwidth}
  \centering
  \includegraphics[width=\textwidth]{sens_conv_background_template}
  \caption{2D background template $T^{\mathrm{bkg}}(\cos\eta,\,E_{\mathrm{reco}})$:
    the 1D spectrum $b(E)$ tiled over nadir and reweighted by $p(\cos\eta)$.}
  \label{fig:sens_conv_background_template}
\end{minipage}\hfill
\begin{minipage}{0.48\textwidth}
  \centering
  \includegraphics[width=\textwidth]{sens_conv_signal_template}
  \caption{2D signal template $T^{\mathrm{sig}}(\cos\eta,\,E_{\mathrm{reco}})$:
    result of convolving the smearing matrix $h$ with the oscillogram
    $P(\nu_e\!\to\!\nu_e)$.  The nadir-dependent wiggles encode the MSW
    day-night modulation exploited by the log-likelihood ratio.}
  \label{fig:sens_conv_signal_template}
\end{minipage}
\end{figure}

\subsection{Results}
\label{sec:sens_results}

Figures~\ref{fig:sens_central_solar_sin12}--\ref{fig:sens_lateral_significance}
show the Sensitivity analysis outputs for both HD configurations
(SolarEnergy, nominal, $30\,\mathrm{kt}\!\cdot\!\mathrm{yr}$).
The top row presents $\sin^2\theta_{12}$ contours for the solar and reactor
hypotheses; the bottom row shows significance vs.\ exposure.

\begin{figure}[ht]
\centering
\begin{minipage}{0.48\textwidth}
  \centering
  \includegraphics[width=\textwidth]{sens_central_solar_sin12}
  \caption{centralAPA: $\Delta\chi^2$ contours in $(\sin^2\theta_{12},\,\Delta m^2_{21})$ under the solar oscillation hypothesis.}
  \label{fig:sens_central_solar_sin12}
\end{minipage}\hfill
\begin{minipage}{0.48\textwidth}
  \centering
  \includegraphics[width=\textwidth]{sens_central_react_sin12}
  \caption{centralAPA: $\Delta\chi^2$ contours under the reactor oscillation hypothesis.}
  \label{fig:sens_central_react_sin12}
\end{minipage}
\end{figure}

\begin{figure}[ht]
\centering
\begin{minipage}{0.48\textwidth}
  \centering
  \includegraphics[width=\textwidth]{sens_central_significance}
  \caption{centralAPA: hypothesis-separation significance vs.\ exposure at 30\,kt$\cdot$yr.}
  \label{fig:sens_central_significance}
\end{minipage}\hfill
\begin{minipage}{0.48\textwidth}
  \centering
  \includegraphics[width=\textwidth]{sens_lateral_solar_sin12}
  \caption{lateralAPA: $\Delta\chi^2$ contours under the solar hypothesis.}
  \label{fig:sens_lateral_solar_sin12}
\end{minipage}
\end{figure}

\begin{figure}[ht]
\centering
\begin{minipage}{0.48\textwidth}
  \centering
  \includegraphics[width=\textwidth]{sens_lateral_react_sin12}
  \caption{lateralAPA: $\Delta\chi^2$ contours under the reactor hypothesis.}
  \label{fig:sens_lateral_react_sin12}
\end{minipage}\hfill
\begin{minipage}{0.48\textwidth}
  \centering
  \includegraphics[width=\textwidth]{sens_lateral_significance}
  \caption{lateralAPA: hypothesis-separation significance vs.\ exposure.}
  \label{fig:sens_lateral_significance}
\end{minipage}
\end{figure}

% ===========================================================================
\section{Comparison Across Analyses}
\label{sec:comparison}

Table~\ref{tab:comparison} summarizes the statistical model for each
analysis.

\begin{table}[ht]
\centering
\caption{Statistical model comparison: Day-Night, HEP, and Sensitivity.}
\label{tab:comparison}
\begin{tabular}{p{3.2cm}p{3.5cm}p{3.5cm}p{3.5cm}}
\toprule
Feature & Day-Night & HEP & Sensitivity \\
\midrule
Observable
  & 1D rate-difference spectrum
  & 1D energy spectrum
  & 2D (energy $\times$ azimuth) \\
Signal
  & $\si{=}\Ecal\xi(r_i^{\rm night}{-}r_i^{\rm day})$
  & $\si{=}\Ecal r_i^{\rm hep}$
  & $\Tsig(\vec\theta)$ \\
Goal
  & Detect day-night asymmetry
  & Discover hep flux
  & Map $(\Delta m^2,\sin^2\theta)$ contour \\
Significance type
  & Gaussian, two-sample Asimov
  & Gaussian, Asimov, PL
  & $\chi^2$ surface \\
Nuisances
  & None (error propagated)
  & $1\times\beta$ (background)
  & $\Apred + \Abkg$ \\
Nuisance solution
  & ---
  & Closed-form quadratic (Eq.~\ref{eq:beta_solution})
  & L-BFGS-B (2D) \\
Penalty form
  & ---
  & $[(\bhat{-}1)/\srel]^2$
  & $(\Apred/\spred)^2{+}(\Abkg/\sbkg)^2$ \\
Per-bin deviance
  & Eq.~\eqref{eq:dn_llr_bin}
  & Eq.~\eqref{eq:q0_expanded} per bin
  & Eq.~\eqref{eq:sens_deviance} \\
Effective background
  & $\Beff$ (Eq.~\ref{eq:dn_Beff})
  & $\bi = \Ecal\sum_{c\neq\mathrm{hep}} r_i^c$
  & $\Tbkg$ (2D template) \\
Barlow-Beeston mask
  & No
  & $N_i^{\rm MC}\ge N_{\rm MC}^{\min}$
  & $\Tbkg>0$ \\
Asymmetry/signal band
  & $\xi\in\{1\pm\epstot,1\}$
  & $d\in\{2.9,3.0,3.1\}$
  & --- \\
PL post-processing
  & ---
  & Gauss($\sigma{=}6$) + PAVA (optional)
  & --- \\
Spike filter
  & ---
  & $\Delta_{\max}$ on pre-PAVA
  & --- \\
Output
  & $Z$ vs.\ exposure
  & $Z$ vs.\ exposure
  & $\chi^2$ map over grid \\
\bottomrule
\end{tabular}
\end{table}

\medskip
\noindent\textbf{Narrative progression.}
The three analyses exhibit a clear progression in statistical complexity
centred on the treatment of nuisance parameters.

The Day-Night analysis makes no nuisance-parameter approximations:
signal and background are evaluated exactly in two separate exposure
intervals, and the asymmetry uncertainty is propagated analytically as a
scale factor on the signal.
The three-source background uncertainty in the Error Gaussian
(Eq.~\eqref{eq:dn_sigma_period}) is the most detailed error propagation
of the three analyses, but no LLR profiling is performed.

The HEP analysis introduces the full profile-likelihood framework.
The critical observation is that a single global $\beta$ correlated across
all bins leads to a scalar stationarity equation, whose positive root is
a closed-form quadratic.  This makes the PL computation exact and fast.
PAVA post-processing and adaptive rebinning address practical numerical
issues (oscillations at low $S/B$; low-statistics tail bins).

The Sensitivity analysis takes the same Baker-Cousins Poisson deviance
as HEP's per-bin LLR, but extends it to 2D and introduces a second
nuisance $\Apred$ that couples the model both to signal and background
simultaneously.  This coupling breaks the factorization that enabled the
HEP quadratic, requiring a full numerical 2D minimization.
The result is a $\chi^2$ map rather than a significance curve.

% ===========================================================================
\section{Summary of Significance Outputs}
\label{sec:summary}

Table~\ref{tab:methods} lists all significance output columns.

\begin{table}[ht]
\centering
\caption{Significance output columns across analyses.}
\label{tab:methods}
\begin{tabular}{llllll}
\toprule
Analysis & Label & Histogram & Rebinning & Estimator & Monotone? \\
\midrule
\multicolumn{6}{l}{\textit{Day-Night}} \\
 & \texttt{RawGaussian}              & Raw    & None & Gaussian         & No \\
 & \texttt{RawErrorGaussian}         & Raw    & None & Gaussian+$\seff$ & No \\
 & \texttt{Gaussian}                 & Smooth & None & Gaussian         & No \\
 & \texttt{ErrorGaussian}            & Smooth & None & Gaussian+$\seff$ & No \\
 & \texttt{RawAsimov}                & Raw    & None & Two-sample LLR   & No \\
 & \texttt{Asimov}                   & Smooth & None & Two-sample LLR   & No \\
\midrule
\multicolumn{6}{l}{\textit{HEP}} \\
 & \texttt{RawGaussianNoRebin}       & Raw    & None     & Gaussian & No \\
 & \texttt{RawAsimovNoRebin}         & Raw    & None     & Asimov   & No \\
 & \texttt{RawGaussian}              & Raw    & Adaptive & Gaussian & Yes \\
 & \texttt{RawAsimov}                & Raw    & Adaptive & Asimov   & Yes \\
 & \texttt{GaussianNoRebin}          & Smooth & None     & Gaussian & No \\
 & \texttt{AsimovNoRebin}            & Smooth & None     & Asimov   & No \\
 & \texttt{Gaussian}                 & Smooth & Adaptive & Gaussian & Yes \\
 & \texttt{Asimov}                   & Smooth & Adaptive & Asimov   & Yes \\
 & \texttt{RawProfileLikelihood}     & Raw    & None     & PL       & Yes (post) \\
 & \texttt{ProfileLikelihood}        & Smooth & None     & PL       & Yes (post) \\
 & \texttt{PreIsotonicProfileLikelihood}
   & Smooth & None & PL (pre-PAVA) & No \\
\midrule
\multicolumn{6}{l}{\textit{Sensitivity}} \\
 & \texttt{chi2\_solar/react} & 2D smooth & None & Baker-Cousins & --- \\
\bottomrule
\end{tabular}
\end{table}

% ===========================================================================
\section{Workflow Flags and Configuration}
\label{sec:flags}

\begin{table}[ht]
\centering
\caption{Orchestrator flags in \texttt{10SensitivityAnalysis.py}.}
\label{tab:flags}
\begin{tabular}{llp{8cm}}
\toprule
Flag & Default & Effect when disabled \\
\midrule
\texttt{--computation}     & True  & Skip all computation; run only plot macros \\
\texttt{--significance}    & True  & Skip \texttt{12DayNight.py}, \texttt{13HEP.py}, \texttt{14Sensitivity*.py} \\
\texttt{--fiducialization} & True  & Skip \texttt{0XFiducializeSignal.py}, \texttt{0YBestFiducial.py} \\
\texttt{--rebin}           & True  & Skip \texttt{11AnalysisSignal.py} adaptive rebinning \\
\texttt{--optimization}    & False & Run smoothing sigma optimizer \texttt{optimize\_smoothing.py} \\
\bottomrule
\end{tabular}
\end{table}

\medskip\noindent
Per-analysis workflow feature flags are controlled by
\texttt{analysis/config.json} under \texttt{WORKFLOW.\{ANALYSIS\}}:

\begin{table}[ht]
\centering
\caption{Per-analysis workflow feature flags.}
\label{tab:analysis_flags}
\begin{tabular}{llp{8cm}}
\toprule
Key & Default & Effect when true \\
\midrule
\texttt{DAYNIGHT.background\_error} & \texttt{true} &
  Compute Error Gaussian curves (Eq.~\ref{eq:dn_error_gaussian}) \\
\texttt{DAYNIGHT.significance\_bins} & \texttt{true} &
  Save per-bin significance spectra at the display exposure \\
\texttt{HEP.pl\_signal\_bands} & \texttt{true} &
  Evaluate three signal normalizations for $\pm1\sigma_s$ PL bands \\
\texttt{HEP.pl\_isotonic} & \texttt{false} &
  Apply Gaussian+PAVA post-processing to PL curves \\
\texttt{HEP.significance\_bins} & \texttt{false} &
  Save per-bin significance spectra at the display exposure \\
\bottomrule
\end{tabular}
\end{table}

% ===========================================================================
\section*{References}
\addcontentsline{toc}{section}{References}
\begin{thebibliography}{9}

\bibitem{Cowan2010}
G.~Cowan, K.~Cranmer, E.~Gross, O.~Vitells,
\textit{Asymptotic formulae for likelihood-based tests of new physics},
Eur.\ Phys.\ J.\ C \textbf{71} (2011) 1554.
\url{https://arxiv.org/abs/1007.1727}

\bibitem{BakerCousins1984}
S.~Baker, R.D.~Cousins,
\textit{Clarification of the use of chi-square and likelihood functions in fits to histograms},
Nucl.\ Instrum.\ Meth.\ \textbf{221} (1984) 437--442.
\url{https://doi.org/10.1016/0167-5087(84)90016-4}

\bibitem{BarlowBeeston1993}
R.~Barlow, C.~Beeston,
\textit{Fitting using finite Monte Carlo samples},
Comput.\ Phys.\ Commun.\ \textbf{77} (1993) 219--228.
\url{https://doi.org/10.1016/0010-4655(93)90005-W}

\bibitem{Robertson1988}
T.~Robertson, F.T.~Wright, R.L.~Dykstra,
\textit{Order Restricted Statistical Inference},
Wiley, 1988.

\bibitem{ByrdEtAl1995}
R.H.~Byrd, P.~Lu, J.~Nocedal, C.~Zhu,
\textit{A limited memory algorithm for bound constrained optimization},
SIAM J.\ Sci.\ Comput.\ \textbf{16} (1995) 1190--1208.

\bibitem{ZhuEtAl1997}
C.~Zhu, R.H.~Byrd, P.~Lu, J.~Nocedal,
\textit{Algorithm 778: L-BFGS-B: Fortran subroutines for large-scale bound-constrained optimization},
ACM Trans.\ Math.\ Softw.\ \textbf{23} (1997) 550--560.

\end{thebibliography}

\end{document}
